\begin{helper}
Fix a finite terminal coordinate set \(U\), a nodup ordered RISI scan list
\(o\) with \(\operatorname{toFinset}(o)=U\), a realized terminal truth
table \(W\subseteq U\), a branch label \(\beta\), and a transcript
\(\tau=(P_\tau,A_\tau,Z_\tau)\).  Let \(F_\tau\subseteq P_\tau\) be the
fixed present support.  Assume
\[
P_\tau,A_\tau\subseteq U,\qquad
Z_\tau\subseteq A_\tau,\qquad
P_\tau\cap A_\tau=\emptyset,
\]
and assume every coordinate of \(P_\tau\) is present in \(W\), while every
coordinate of \(A_\tau\) is absent in \(W\).

For each terminal coordinate \(e\), let \(D_\tau(e)\) be the set of terminal
coordinates read by the RISI scan when it evaluates, at \(e\),
staged-openness, touchedness, same-color closedness, and pre-terminal
recordedness.  The generated scan interface consists of the following
explicit clauses.

First, every read coordinate has already been classified by the transcript:
for every \(e\in U\) and every \(c\in D_\tau(e)\), the coordinate \(c\) is
one of
\[
F_\tau,\qquad P_\tau\setminus F_\tau,\qquad
Z_\tau,\qquad A_\tau\setminus Z_\tau .
\]
Moreover, if \(o=\mathrm{pre}\, e\,\mathrm{suffix}\), then
\[
D_\tau(e)\subseteq \operatorname{toFinset}(\mathrm{pre});
\]
that is, the scan reads only earlier terminal coordinates.  In the intended
branch/transcript instantiation, this past-read condition is obtained by
applying the terminal truth-table extension
\(\noderef{FixedSetHistoryCellTerminalAssignmentExtension}\) and then using
the prefix-safety hypothesis of the branch/transcript cell.

Second, the sets \(D_\tau(e)\) are exact read sets for the four tests: for
any terminal truth table \(A\subseteq U\), if \(A\) and \(W\) agree on
\(D_\tau(e)\), then the staged-open, touched, same-color-closed, and
pre-terminal-recorded tests at \(e\) have the same values for \(A\) and
\(W\).  The selected absent coordinates used in this generated interface
are the canonical selected absences supplied by
\(\noderef{FixedSetHistoryCellCanonicalAbsenceSelection}\).

Third, the RISI scan is deterministic from the ordered list and these four
test outcomes: if a terminal truth table \(A\subseteq U\) has the same four
test outcomes as \(W\) at every \(e\in U\), then the scan returns the same
branch label and transcript,
\[
\operatorname{branchOfAssignment}(A)=\beta,\qquad
\operatorname{scanTranscript}(A)=\operatorname{some}(\tau).
\]

Then the generated interface supplies exactly the three consequences needed
by the residual scan-stability step:
\[
D_\tau(e)\subseteq P_\tau\cup A_\tau\qquad(e\in U),
\]
agreement on \(D_\tau(e)\) is enough to preserve all four tests at \(e\),
and equality of all four test outcomes over \(U\) forces scan output
\((\beta,\tau)\).
\end{helper}
\begin{proof}
Fix \(e\in U\) and \(c\in D_\tau(e)\).  By the generated classification
clause there are four cases.

If \(c\in F_\tau\), then \(F_\tau\subseteq P_\tau\), so
\(c\in P_\tau\subseteq P_\tau\cup A_\tau\).  If
\(c\in P_\tau\setminus F_\tau\), then again \(c\in P_\tau\), hence
\(c\in P_\tau\cup A_\tau\).  If \(c\in Z_\tau\), then
\(Z_\tau\subseteq A_\tau\), so \(c\in A_\tau\subseteq P_\tau\cup A_\tau\).
Finally, if \(c\in A_\tau\setminus Z_\tau\), then \(c\in A_\tau\), hence
\(c\in P_\tau\cup A_\tau\).  Since \(c\) was arbitrary, this proves
\[
D_\tau(e)\subseteq P_\tau\cup A_\tau .
\]

The same classification also explains why the RISI read is a legal prefix
read in the intended scan: whenever \(o=\mathrm{pre}\,e\,\mathrm{suffix}\),
the generated past-read clause gives
\(D_\tau(e)\subseteq\operatorname{toFinset}(\mathrm{pre})\).  Thus the
branch/transcript prefix-safety hypothesis applies precisely to the
coordinates in \(D_\tau(e)\); no later terminal coordinate is used in the
test at \(e\).

Now let \(A\subseteq U\) be a terminal truth table which agrees with
\(W\) on \(D_\tau(e)\).  The exact-read-stability clause says directly
that this agreement preserves, at the coordinate \(e\), each of the four
quantities read by the scan: staged-openness, touchedness, same-color
closedness, and pre-terminal recordedness.  This proves the second exported
consequence.

Finally suppose \(A\subseteq U\) has the same four test outcomes as \(W\) at
every terminal coordinate.  The deterministic-output clause of the generated
RISI scan interface applies to this whole test sequence.  Therefore the
branch chosen from \(A\) is \(\beta\), and the transcript returned by the
scan on \(A\) is \(\operatorname{some}(\tau)\).  This is the third exported
consequence.
\end{proof}
